% !TeX program = xelatex
% !TeX encoding = UTF-8
%%%%%%%%%%%%%%%%%%%%%%%%%%%%%%%%%%%%%%%%%%%%%%%%%%%%%%%%%%%%%%%%%%%%%%%%%%%%%%%%
%%
%%   قالب اسلایدهای ارائهٔ دفاع پایان‌نامه  —  PersianThesis beamer theme
%%
%%   کامپایل:   xelatex main.tex      (دو بار، برای فهرست مطالب و شمارهٔ اسلایدها)
%%   یا:        latexmk main.tex
%%   یا:        make
%%
%%   ساختار پروژه
%%   -------------
%%     main.tex               همین فایل: چیدمان کلی و ترتیب فصل‌ها
%%     config/packages.tex    بسته‌های عمومی
%%     config/info.tex        ★ مشخصات دفاع؛ معمولاً تنها فایلی که باید ویرایش شود
%%     config/fonts.tex       فونت‌ها و xepersian  (همیشه آخرین ورودی پیش‌درآمد)
%%     universities/*.tex     پیش‌تنظیم رنگ و لوگوی دانشگاه‌ها
%%     content/*.tex          محتوای اسلایدها، هر فصل یک فایل
%%     theme/                 خود قالب؛ برای استفادهٔ معمولی نیازی به تغییر ندارد
%%     figures/, refs/        تصاویر و منابع
%%
%%   نکتهٔ مهم دربارهٔ ترتیب پیش‌درآمد
%%   ------------------------------
%%   بستهٔ xepersian باید آخرین بستهٔ پیش‌درآمد باشد؛ بنابراین
%%   config/fonts.tex همیشه در پایان و بلافاصله پیش از \begin{document} می‌آید.
%%
%%%%%%%%%%%%%%%%%%%%%%%%%%%%%%%%%%%%%%%%%%%%%%%%%%%%%%%%%%%%%%%%%%%%%%%%%%%%%%%%
\documentclass[11pt,aspectratio=169]{beamer}

%%%%%%%%%%%%%%%%%%%%%%%%%%%%%%%%%%%%%%%%%%%%%%%%%%%%%%%%%%%%%%%%%%%%%%%%%%%%%%%%
%% ۰. مسیرها  —  تا LaTeX قالب و محتوا را در زیرپوشه‌ها پیدا کند
%%%%%%%%%%%%%%%%%%%%%%%%%%%%%%%%%%%%%%%%%%%%%%%%%%%%%%%%%%%%%%%%%%%%%%%%%%%%%%%%
\makeatletter
\def\input@path{{./}{theme/}{config/}{content/}{universities/}}
\makeatother

%%%%%%%%%%%%%%%%%%%%%%%%%%%%%%%%%%%%%%%%%%%%%%%%%%%%%%%%%%%%%%%%%%%%%%%%%%%%%%%%
%% ۱. بسته‌های عمومی
%%%%%%%%%%%%%%%%%%%%%%%%%%%%%%%%%%%%%%%%%%%%%%%%%%%%%%%%%%%%%%%%%%%%%%%%%%%%%%%%
%%%%%%%%%%%%%%%%%%%%%%%%%%%%%%%%%%%%%%%%%%%%%%%%%%%%%%%%%%%%%%%%%%%%%%%%%%%%%%%%
%%  config/packages.tex
%%  Generic LaTeX packages used by the slides.
%%
%%  Deliberately short.  Everything here ships with a normal TeX Live / MiKTeX
%%  installation, and nothing here touches fonts or text direction -- those live
%%  in config/fonts.tex, which must stay the *last* thing loaded in the preamble
%%  because of xepersian.
%%%%%%%%%%%%%%%%%%%%%%%%%%%%%%%%%%%%%%%%%%%%%%%%%%%%%%%%%%%%%%%%%%%%%%%%%%%%%%%%

% ---- mathematics -----------------------------------------------------------
\usepackage{amsmath}
\usepackage{amssymb}

% ---- tables ----------------------------------------------------------------
\usepackage{array}
\usepackage{booktabs}
\usepackage{multirow}
\usepackage{colortbl}      % \rowcolor / \cellcolor

% ---- graphics --------------------------------------------------------------
% The theme pulls these in as well; loading them here too costs nothing and
% keeps this file usable on its own.
\usepackage{graphicx}
\usepackage{tikz}
\usetikzlibrary{positioning}
\usetikzlibrary{arrows.meta}
\usetikzlibrary{calc}
\usetikzlibrary{fit}
\usetikzlibrary{shapes.geometric}
\usetikzlibrary{backgrounds}

% ---- optional extras ------------------------------------------------------
% Uncomment if you need real statistical plots.  The template itself draws its
% charts in plain TikZ, so pgfplots is not required for it to compile.
%\usepackage{pgfplots}
%\pgfplotsset{compat=newest}

% Table cells that keep their own direction, handy in mixed Persian/English
% tables:  \begin{tabular}{...} ... \LTRfootnote / \lr{...}
% (\lr comes from bidi and is available everywhere in the document body.)


\graphicspath{{figures/}{figures/logos/}}

%%%%%%%%%%%%%%%%%%%%%%%%%%%%%%%%%%%%%%%%%%%%%%%%%%%%%%%%%%%%%%%%%%%%%%%%%%%%%%%%
%% ۲. قالب
%%
%%  همهٔ کلیدها در theme/beamerthemePersianThesis.sty فهرست شده‌اند.
%%  palette: indigo | navy | teal | maroon | forest | plum | graphite
%%  mode:    light | dark
%%%%%%%%%%%%%%%%%%%%%%%%%%%%%%%%%%%%%%%%%%%%%%%%%%%%%%%%%%%%%%%%%%%%%%%%%%%%%%%%
\usetheme[palette=indigo,mode=light]{PersianThesis}

\PTsetup{
  progressbar = foot,        % none | foot | head | both
  numbering   = fraction,    % fraction | plain | none
  footline    = true,
  eyebrow     = true,        % «فصل ۲ : ...» بالای عنوان هر اسلاید
  sectionpage = true,        % اسلاید جداکنندهٔ ابتدای هر فصل
  tocnumbers  = true,
  margin      = 1.05cm,
}

%%%%%%%%%%%%%%%%%%%%%%%%%%%%%%%%%%%%%%%%%%%%%%%%%%%%%%%%%%%%%%%%%%%%%%%%%%%%%%%%
%% ۳. پیش‌تنظیم دانشگاه  (اختیاری)
%%
%%  رنگ‌ها، لوگو و نام دانشگاه را یک‌جا تنظیم می‌کند.  هر چیزی که بعد از آن در
%%  config/info.tex نوشته شود، بر این مقادیر مقدم است.
%%%%%%%%%%%%%%%%%%%%%%%%%%%%%%%%%%%%%%%%%%%%%%%%%%%%%%%%%%%%%%%%%%%%%%%%%%%%%%%%
%%%%%%%%%%%%%%%%%%%%%%%%%%%%%%%%%%%%%%%%%%%%%%%%%%%%%%%%%%%%%%%%%%%%%%%%%%%%%%%%
%%  universities/sharif.tex — دانشگاه صنعتی شریف
%%
%%  The colours are a careful visual match to the university's blue, not an
%%  official brand specification.  If your faculty publishes exact values, put
%%  them in the primary/secondary/accent keys below.
%%%%%%%%%%%%%%%%%%%%%%%%%%%%%%%%%%%%%%%%%%%%%%%%%%%%%%%%%%%%%%%%%%%%%%%%%%%%%%%%

\PTsetup{
  primary   = 1C3F72,
  secondary = 33639E,
  accent    = C08A2E,
  logo      = figures/logos/sharif.pdf,
}

\university{دانشگاه صنعتی شریف}
\shortuniversity{دانشگاه صنعتی شریف}
\faculty{دانشکدهٔ مهندسی کامپیوتر}
\department{}


%%%%%%%%%%%%%%%%%%%%%%%%%%%%%%%%%%%%%%%%%%%%%%%%%%%%%%%%%%%%%%%%%%%%%%%%%%%%%%%%
%% ۴. مشخصات دفاع
%%%%%%%%%%%%%%%%%%%%%%%%%%%%%%%%%%%%%%%%%%%%%%%%%%%%%%%%%%%%%%%%%%%%%%%%%%%%%%%%
%%%%%%%%%%%%%%%%%%%%%%%%%%%%%%%%%%%%%%%%%%%%%%%%%%%%%%%%%%%%%%%%%%%%%%%%%%%%%%%%
%%  config/info.tex
%%  ★ THIS IS THE ONLY FILE MOST USERS NEED TO EDIT ★
%%
%%  All the metadata of the defence.  Leave a field empty (or delete the line)
%%  and the corresponding piece of the title page disappears by itself -- a
%%  bachelor project with a single supervisor and no referee still produces a
%%  balanced, finished-looking cover.
%%%%%%%%%%%%%%%%%%%%%%%%%%%%%%%%%%%%%%%%%%%%%%%%%%%%%%%%%%%%%%%%%%%%%%%%%%%%%%%%

%%%%%%%%%%%%%%%%%%%%%%%%%%%%%%%%%%%%%%%%%%%%%%%%%%%%%%%%%%%%%%%%%%%%%%%%%%%%%%%%
%% 1. Institution
%%%%%%%%%%%%%%%%%%%%%%%%%%%%%%%%%%%%%%%%%%%%%%%%%%%%%%%%%%%%%%%%%%%%%%%%%%%%%%%%
\university{دانشگاه صنعتی شریف}
\shortuniversity{دانشگاه صنعتی شریف}   % used in the footer; empty = full name
\faculty{دانشکدهٔ مهندسی کامپیوتر}
\department{}                            % گروه/رشته, if your faculty needs it

%%%%%%%%%%%%%%%%%%%%%%%%%%%%%%%%%%%%%%%%%%%%%%%%%%%%%%%%%%%%%%%%%%%%%%%%%%%%%%%%
%% 2. Degree
%%
%%  The cover line reads:  <thesisfor label> <degree> <major>
%%  For a bachelor's final project the usual wording is a different label:
%%      \PTsetlabel{thesisfor}{پروژهٔ پایانی دورهٔ}
%%%%%%%%%%%%%%%%%%%%%%%%%%%%%%%%%%%%%%%%%%%%%%%%%%%%%%%%%%%%%%%%%%%%%%%%%%%%%%%%
\degree{کارشناسی}                       % کارشناسی / کارشناسی ارشد / دکتری
\major{مهندسی کامپیوتر}                 % گرایش, optional
\academicyear{۱۴۰۴--۱۴۰۵}

%%%%%%%%%%%%%%%%%%%%%%%%%%%%%%%%%%%%%%%%%%%%%%%%%%%%%%%%%%%%%%%%%%%%%%%%%%%%%%%%
%% 3. Title
%%
%%  \title      Persian title, the largest thing on the cover
%%  \latintitle English title, set in the Latin font one size down (optional)
%%  \subtitle   optional one-liner underneath
%%  \shorttitle goes in the middle cell of the footer, so keep it short
%%%%%%%%%%%%%%%%%%%%%%%%%%%%%%%%%%%%%%%%%%%%%%%%%%%%%%%%%%%%%%%%%%%%%%%%%%%%%%%%
\title{سامانهٔ پرسش و پاسخ هوشمند مبتنی بر وب}
\latintitle{Intelligent Web-Based Question Answering System}
\subtitle{}
\shorttitle{سامانهٔ پرسش و پاسخ هوشمند}

%%%%%%%%%%%%%%%%%%%%%%%%%%%%%%%%%%%%%%%%%%%%%%%%%%%%%%%%%%%%%%%%%%%%%%%%%%%%%%%%
%% 4. People
%%
%%  Only \author is mandatory.  Every other name may be left out.
%%%%%%%%%%%%%%%%%%%%%%%%%%%%%%%%%%%%%%%%%%%%%%%%%%%%%%%%%%%%%%%%%%%%%%%%%%%%%%%%
\author{علی هاشمیان}
\studentid{}                             % شمارهٔ دانشجویی, optional
\supervisor{دکتر حمید بیگی}
\supervisortwo{}
\advisor{}
\advisortwo{}
\referee{}
\refereetwo{}

%%%%%%%%%%%%%%%%%%%%%%%%%%%%%%%%%%%%%%%%%%%%%%%%%%%%%%%%%%%%%%%%%%%%%%%%%%%%%%%%
%% 5. Date and place of the defence
%%%%%%%%%%%%%%%%%%%%%%%%%%%%%%%%%%%%%%%%%%%%%%%%%%%%%%%%%%%%%%%%%%%%%%%%%%%%%%%%
\date{مرداد ۱۴۰۵}
\defenseplace{}                          % e.g. سالن سمینار دانشکده

%%%%%%%%%%%%%%%%%%%%%%%%%%%%%%%%%%%%%%%%%%%%%%%%%%%%%%%%%%%%%%%%%%%%%%%%%%%%%%%%
%% 6. Invocation
%%
%%  Printed in small type above the logo.  Leave empty to omit it entirely.
%%%%%%%%%%%%%%%%%%%%%%%%%%%%%%%%%%%%%%%%%%%%%%%%%%%%%%%%%%%%%%%%%%%%%%%%%%%%%%%%
\invocation{بسم الله الرحمن الرحیم}

%%%%%%%%%%%%%%%%%%%%%%%%%%%%%%%%%%%%%%%%%%%%%%%%%%%%%%%%%%%%%%%%%%%%%%%%%%%%%%%%
%% 7. Wording
%%
%%  Every fixed string the theme prints can be replaced here.  The full list of
%%  keys is documented at the top of theme/persian-thesis-utils.sty.
%%%%%%%%%%%%%%%%%%%%%%%%%%%%%%%%%%%%%%%%%%%%%%%%%%%%%%%%%%%%%%%%%%%%%%%%%%%%%%%%
%\PTsetlabel{thesisfor}{پروژهٔ پایانی دورهٔ}
%\PTsetlabel{student}{ارائه‌دهنده}
%\PTsetlabel{thanks}{با تشکر از توجه شما}
%\PTsetlabel{chapter}{بخش}

%%%%%%%%%%%%%%%%%%%%%%%%%%%%%%%%%%%%%%%%%%%%%%%%%%%%%%%%%%%%%%%%%%%%%%%%%%%%%%%%
%% 8. PDF metadata
%%%%%%%%%%%%%%%%%%%%%%%%%%%%%%%%%%%%%%%%%%%%%%%%%%%%%%%%%%%%%%%%%%%%%%%%%%%%%%%%
\hypersetup{
  pdftitle    = {سامانهٔ پرسش و پاسخ هوشمند مبتنی بر وب},
  pdfauthor   = {علی هاشمیان},
  pdfsubject  = {ارائهٔ دفاع پایان‌نامهٔ کارشناسی},
  pdfkeywords = {پرسش و پاسخ, بازیابی اطلاعات, مدل زبانی بزرگ, RAG},
  colorlinks  = true,
  linkcolor   = PTprimary,
  urlcolor    = PTsecondary,
  citecolor   = PTaccentdark,
}


%%%%%%%%%%%%%%%%%%%%%%%%%%%%%%%%%%%%%%%%%%%%%%%%%%%%%%%%%%%%%%%%%%%%%%%%%%%%%%%%
%% ۵. فونت‌ها و xepersian  —  همیشه آخرین بخش پیش‌درآمد
%%%%%%%%%%%%%%%%%%%%%%%%%%%%%%%%%%%%%%%%%%%%%%%%%%%%%%%%%%%%%%%%%%%%%%%%%%%%%%%%
%%%%%%%%%%%%%%%%%%%%%%%%%%%%%%%%%%%%%%%%%%%%%%%%%%%%%%%%%%%%%%%%%%%%%%%%%%%%%%%%
%%  config/fonts.tex
%%  Fonts and Persian typesetting.  **This file must be the last thing loaded in
%%  the preamble**, because xepersian (and the bidi package behind it) has to be
%%  the last package of the document.
%%
%%  The bundled family is Vazirmatn by Saber Rastikerdar (SIL OFL 1.1), in
%%  fonts/Vazirmatn/.  It covers Persian *and* Latin, so one family carries the
%%  whole deck and mixed Persian/English lines stay optically consistent.
%%
%%  Paths are relative to the directory you compile from, i.e. the folder that
%%  holds main.tex.  To use another family, change the four \...font commands
%%  below and nothing else -- the theme never refers to a concrete font file.
%%%%%%%%%%%%%%%%%%%%%%%%%%%%%%%%%%%%%%%%%%%%%%%%%%%%%%%%%%%%%%%%%%%%%%%%%%%%%%%%

% Proper Computer Modern maths (Vazirmatn has no maths glyphs).  Loaded *before*
% xepersian on purpose: no package may come after it.
\usefonttheme[onlymath]{serif}

%%%%%%%%%%%%%%%%%%%%%%%%%%%%%%%%%%%%%%%%%%%%%%%%%%%%%%%%%%%%%%%%%%%%%%%%%%%%%%%%
%% xepersian -- the last package in the preamble
%%%%%%%%%%%%%%%%%%%%%%%%%%%%%%%%%%%%%%%%%%%%%%%%%%%%%%%%%%%%%%%%%%%%%%%%%%%%%%%%
\usepackage{xepersian}

%%%%%%%%%%%%%%%%%%%%%%%%%%%%%%%%%%%%%%%%%%%%%%%%%%%%%%%%%%%%%%%%%%%%%%%%%%%%%%%%
%% Main Persian font
%%
%%  Mapping=parsidigits turns every ASCII digit of the Persian text into a
%%  Persian one, so slide numbers, chapter numbers and outline numbers come out
%%  as ۱ ۲ ۳ without any extra markup.
%%  Persian has no italics; the italic slot is mapped to the Medium weight, which
%%  gives \emph{...} a gentle emphasis instead of a fake slant.
%%%%%%%%%%%%%%%%%%%%%%%%%%%%%%%%%%%%%%%%%%%%%%%%%%%%%%%%%%%%%%%%%%%%%%%%%%%%%%%%
\settextfont[
  Path            = fonts/Vazirmatn/,
  Extension       = .ttf,
  UprightFont     = *-Regular,
  BoldFont        = *-Bold,
  ItalicFont      = *-Medium,
  BoldItalicFont  = *-ExtraBold,
  Mapping         = parsidigits,
  Scale           = 1.0,
]{Vazirmatn}

%%%%%%%%%%%%%%%%%%%%%%%%%%%%%%%%%%%%%%%%%%%%%%%%%%%%%%%%%%%%%%%%%%%%%%%%%%%%%%%%
%% Latin font
%%
%%  Same family, but *without* parsidigits: English runs keep Latin digits, so
%%  "Precision@3" or "GPT-4" look the way they should.
%%%%%%%%%%%%%%%%%%%%%%%%%%%%%%%%%%%%%%%%%%%%%%%%%%%%%%%%%%%%%%%%%%%%%%%%%%%%%%%%
\setlatintextfont[
  Path            = fonts/Vazirmatn/,
  Extension       = .ttf,
  UprightFont     = *-Regular,
  BoldFont        = *-Bold,
  ItalicFont      = *-Medium,
  BoldItalicFont  = *-ExtraBold,
  Ligatures       = TeX,
  Scale           = 1.0,
]{Vazirmatn}

% The family \en{...} switches to.  Keep it in sync with \setlatintextfont.
\newfontfamily\PTlatinfont[
  Path            = fonts/Vazirmatn/,
  Extension       = .ttf,
  UprightFont     = *-Regular,
  BoldFont        = *-Bold,
  ItalicFont      = *-Medium,
  BoldItalicFont  = *-ExtraBold,
  Ligatures       = TeX,
]{Vazirmatn}

% A true semibold, used for \hl{...} and \term{...} instead of a heavy bold.
%
%  Script=Arabic is not optional here.  \settextfont / \setlatintextfont get the
%  Persian defaults from xepersian, but a hand-rolled \newfontfamily does not:
%  without an explicit Arabic script the shaping features never fire and the
%  letters come out unjoined.
\newfontfamily\PTsemibold[
  Path            = fonts/Vazirmatn/,
  Extension       = .ttf,
  Script          = Arabic,
  UprightFont     = *-SemiBold,
  BoldFont        = *-Bold,
  ItalicFont      = *-SemiBold,
  BoldItalicFont  = *-Bold,
  Mapping         = parsidigits,
]{Vazirmatn}

% A light weight for very large display text, should you want it.
\newfontfamily\PTlight[
  Path            = fonts/Vazirmatn/,
  Extension       = .ttf,
  Script          = Arabic,
  UprightFont     = *-Light,
  BoldFont        = *-Medium,
  ItalicFont      = *-Light,
  BoldItalicFont  = *-Medium,
  Mapping         = parsidigits,
]{Vazirmatn}

% Monospace for the PTcode environment.  Left at the LaTeX default (Latin
% Modern Mono), which every distribution has.  Uncomment to upgrade:
%\setlatinmonofont[Scale=MatchLowercase]{DejaVu Sans Mono}

%%%%%%%%%%%%%%%%%%%%%%%%%%%%%%%%%%%%%%%%%%%%%%%%%%%%%%%%%%%%%%%%%%%%%%%%%%%%%%%%
%% Weight tuning for the theme's inline helpers
%%
%%  Now that a real SemiBold exists, \hl{...} and \term{...} can use it: on a
%%  projector a semibold Persian word reads as emphasis, a bold one reads as
%%  shouting.
%%%%%%%%%%%%%%%%%%%%%%%%%%%%%%%%%%%%%%%%%%%%%%%%%%%%%%%%%%%%%%%%%%%%%%%%%%%%%%%%
\setbeamerfont{PT highlight}{family=\PTsemibold,series=\mdseries}
\setbeamerfont{PT term}     {family=\PTsemibold,series=\mdseries}


%%%%%%%%%%%%%%%%%%%%%%%%%%%%%%%%%%%%%%%%%%%%%%%%%%%%%%%%%%%%%%%%%%%%%%%%%%%%%%%%
\begin{document}

% ---- جلد -------------------------------------------------------------------
\PTtitleframe

% ---- فهرست مطالب ----------------------------------------------------------
% \PToutlineframe[hideallsubsections]  حالت پیش‌فرض؛ برای نمایش زیربخش‌ها:
% \PToutlineframe[]
\PToutlineframe

% ---- فصل‌ها ----------------------------------------------------------------
%%%%%%%%%%%%%%%%%%%%%%%%%%%%%%%%%%%%%%%%%%%%%%%%%%%%%%%%%%%%%%%%%%%%%%%%%%%%%%%%
%%  content/01-introduction.tex — فصل ۱: مقدمه و بیان مسئله
%%
%%  هر \section یک اسلاید جداکننده می‌سازد (با sectionpage=false خاموش می‌شود)
%%  و نام همان فصل در نوار بالای عنوان اسلایدها تکرار می‌شود.
%%%%%%%%%%%%%%%%%%%%%%%%%%%%%%%%%%%%%%%%%%%%%%%%%%%%%%%%%%%%%%%%%%%%%%%%%%%%%%%%
\section{مقدمه و بیان مسئله}

%%%%%%%%%%%%%%%%%%%%%%%%%%%%%%%%%%%%%%%%%%%%%%%%%%%%%%%%%%%%%%%%%%%%%%%%%%%%%%%%
\begin{frame}{صورت مسئله}{از فهرست پیوندها تا پاسخ}
  \begin{itemize}
    \item کاربر پرسشی به \term{زبان طبیعی} می‌پرسد، اما موتور جست‌وجو فهرستی از
          پیوندها بازمی‌گرداند؛ یافتن پاسخ در میان آن‌ها بر دوش خود کاربر است.
    \item مدل‌های زبانی بزرگ پاسخ روان می‌دهند، ولی دانش‌شان به \hl{زمان آموزش}
          محدود است و در نتیجه از رخدادها و مستندات تازه بی‌خبرند.
    \item همین محدودیت به \term{توهم‌زایی} می‌انجامد: پاسخی که درست به نظر می‌رسد
          اما پشتوانهٔ قابل بررسی ندارد.
  \end{itemize}

  \begin{PTcallout}
    پرسش پژوهش: چگونه می‌توان \hl{روانی} پاسخ مدل زبانی را با \hl{تازگی} و
    \hl{قابلیت استناد} وب یکی کرد؟
  \end{PTcallout}
\end{frame}

%%%%%%%%%%%%%%%%%%%%%%%%%%%%%%%%%%%%%%%%%%%%%%%%%%%%%%%%%%%%%%%%%%%%%%%%%%%%%%%%
\begin{frame}{کاستی سامانه‌های موجود}
  \begin{itemize}
    \item \term{موتورهای کلیدواژه‌ای}: بازیابی سریع، اما ناتوان از درک منظور
          پرسش‌های چندبخشی.
    \item \term{دستیارهای گفت‌وگویی بستهٔ وب}: پاسخ یکدست، بدون ارجاع بررسی‌پذیر
          به سند اصلی.
    \item \term{سامانه‌های تجاری پرسش و پاسخ}: جعبهٔ سیاه‌اند؛ نه سنجش‌پذیرند و نه
          برای زبان و دامنهٔ خاص قابل تنظیم.
  \end{itemize}

  \begin{PTcons}
    \begin{itemize}
      \item نبود سنجهٔ روشن برای اینکه هر جزء پاسخ از کدام سند آمده است.
      \item هزینهٔ بالای فراخوانی مدل، وقتی همهٔ اسناد بازیابی‌شده به آن سپرده شوند.
    \end{itemize}
  \end{PTcons}
\end{frame}

%%%%%%%%%%%%%%%%%%%%%%%%%%%%%%%%%%%%%%%%%%%%%%%%%%%%%%%%%%%%%%%%%%%%%%%%%%%%%%%%
\begin{frame}{هدف و مشارکت‌های این پژوهش}
  \begin{enumerate}
    \item طراحی سامانه‌ای مبتنی بر وب که پرسش زبان طبیعی را می‌گیرد و
          \hl{پاسخ مستند} با ارجاع به منبع بازمی‌گرداند.
    \item زنجیرهٔ \term{بازیابی، بازرتبه‌بندی و تولید} با اجزای جایگزین‌پذیر، تا
          هر مرحله جداگانه سنجیده شود.
    \item \term{بازرتبه‌بندی دومرحله‌ای} برای کوتاه کردن بافتار ورودی مدل و کاهش
          هزینهٔ هر پرسش.
    \item ارزیابی مقایسه‌ای روی مجموعه‌ای از \hl{۴۰ پرسش} در برابر یک سامانهٔ
          تجاری مرجع.
  \end{enumerate}

  \begin{PTresult}
    دقت پاسخ‌دهی سامانه \hl{۹۵٪} در برابر \hl{۷۷٫۵٪} سامانهٔ مرجع، با بافتاری
    کوتاه‌تر و هزینهٔ کمتر.
  \end{PTresult}
\end{frame}

%%%%%%%%%%%%%%%%%%%%%%%%%%%%%%%%%%%%%%%%%%%%%%%%%%%%%%%%%%%%%%%%%%%%%%%%%%%%%%%%
%%  content/02-background.tex — فصل ۲: مفاهیم پایه و مبانی نظری
%%%%%%%%%%%%%%%%%%%%%%%%%%%%%%%%%%%%%%%%%%%%%%%%%%%%%%%%%%%%%%%%%%%%%%%%%%%%%%%%
\section{مفاهیم پایه و مبانی نظری}

%%%%%%%%%%%%%%%%%%%%%%%%%%%%%%%%%%%%%%%%%%%%%%%%%%%%%%%%%%%%%%%%%%%%%%%%%%%%%%%%
\begin{frame}{بازیابی‌افزودهٔ تولید}{\en{Retrieval-Augmented Generation}}
  \begin{PTdef}[\en{RAG}]
    الگویی که پیش از تولید پاسخ، اسناد مرتبط را از یک منبع بیرونی بازیابی می‌کند
    و آن‌ها را همراه پرسش به مدل زبانی می‌سپارد؛ بنابراین دانش سامانه از
    \hl{وزن‌های مدل} جدا و در \hl{زمان اجرا} قابل به‌روزرسانی است.
  \end{PTdef}

  \begin{itemize}
    \item دانش تازه بدون آموزش دوباره افزوده می‌شود.
    \item هر جملهٔ پاسخ به سند مشخصی گره می‌خورد، پس \term{استنادپذیر} است.
    \item کیفیت پاسخ مستقیماً تابع کیفیت \hl{بازیابی} است، نه فقط اندازهٔ مدل.
  \end{itemize}
\end{frame}

%%%%%%%%%%%%%%%%%%%%%%%%%%%%%%%%%%%%%%%%%%%%%%%%%%%%%%%%%%%%%%%%%%%%%%%%%%%%%%%%
\begin{frame}{بازیابی معنایی}{از تطبیق واژه تا تطبیق معنا}
  پرسش و سند هر دو به برداری در فضایی چندهزار‌بعدی نگاشته می‌شوند و مرتبط‌بودن با
  \term{شباهت کسینوسی} سنجیده می‌شود:
  \[
    \operatorname{sim}(q,d)=
      \frac{\mathbf{v}_q\cdot\mathbf{v}_d}
           {\lVert\mathbf{v}_q\rVert\;\lVert\mathbf{v}_d\rVert}
  \]

  \begin{itemize}
    \item هم‌معنایی و صورت‌های واژگانی گوناگون به‌طور طبیعی پوشش داده می‌شوند.
    \item نگاشت یک‌بار محاسبه و نمایه می‌شود؛ در زمان پاسخ‌دهی فقط جست‌وجوی
          نزدیک‌ترین همسایه لازم است.
  \end{itemize}

  \begin{PTnote}
    بازیابی معنایی \hl{یادآوری} بالایی دارد اما ترتیبش خام است؛ همین شکاف، انگیزهٔ
    مرحلهٔ بازرتبه‌بندی در معماری پیشنهادی است.
  \end{PTnote}
\end{frame}

%%%%%%%%%%%%%%%%%%%%%%%%%%%%%%%%%%%%%%%%%%%%%%%%%%%%%%%%%%%%%%%%%%%%%%%%%%%%%%%%
\begin{frame}{بازرتبه‌بندی}{\en{Cross-Encoder Reranking}}
  \begin{itemize}
    \item بازیاب برداری پرسش و سند را \hl{جداگانه} می‌بیند؛ بازرتبه‌بند آن‌ها را
          \hl{با هم} به مدل می‌دهد و امتیاز مرتبط‌بودن مستقیم می‌گیرد.
    \item هزینهٔ محاسباتی بالاتر است، پس تنها روی \term{چند دهک بالای} نتایج
          بازیابی اجرا می‌شود.
  \end{itemize}

  \begin{PTresult}
    ترکیب «بازیابی گسترده + بازرتبه‌بندی باریک» هم دقت را بالا می‌برد و هم بافتار
    ورودی مدل زبانی را کوتاه می‌کند.
  \end{PTresult}
\end{frame}

%%%%%%%%%%%%%%%%%%%%%%%%%%%%%%%%%%%%%%%%%%%%%%%%%%%%%%%%%%%%%%%%%%%%%%%%%%%%%%%%
\begin{frame}{سنجه‌های ارزیابی}
  \begin{description}
    \item[\en{Precision@k}] نسبت اسناد مرتبط در \(k\) نتیجهٔ نخست.
    \item[\en{NDCG@k}] همان دقت، اما با وزن‌دهی به \hl{جایگاه} هر سند:
  \end{description}
  \[
    \operatorname{DCG}@k=\sum_{i=1}^{k}\frac{2^{r_i}-1}{\log_2(i+1)},
    \qquad
    \operatorname{NDCG}@k=\frac{\operatorname{DCG}@k}{\operatorname{IDCG}@k}
  \]
  \begin{itemize}
    \item \(r_i\) درجهٔ مرتبط‌بودن سند جایگاه \(i\) و \(\operatorname{IDCG}@k\)
          مقدار همان سنجه در ترتیب آرمانی است.
    \item مقدار ۱ یعنی ترتیب بازیابی‌شده با ترتیب آرمانی یکی است.
  \end{itemize}
\end{frame}

%%%%%%%%%%%%%%%%%%%%%%%%%%%%%%%%%%%%%%%%%%%%%%%%%%%%%%%%%%%%%%%%%%%%%%%%%%%%%%%%
%%  content/03-related-work.tex — فصل ۳: کارهای پیشین و تحلیل سامانه‌های موجود
%%%%%%%%%%%%%%%%%%%%%%%%%%%%%%%%%%%%%%%%%%%%%%%%%%%%%%%%%%%%%%%%%%%%%%%%%%%%%%%%
\section{کارهای پیشین و تحلیل سامانه‌های موجود}

%%%%%%%%%%%%%%%%%%%%%%%%%%%%%%%%%%%%%%%%%%%%%%%%%%%%%%%%%%%%%%%%%%%%%%%%%%%%%%%%
\begin{frame}{سه خط پژوهشی}
  \begin{description}
    \item[بازیابی متنی] از تطبیق واژگانی تا بازیاب‌های برداری چگال
          \cite{karpukhin2020dpr}.
    \item[تولید مبتنی بر بازیابی] پیوند دادن بازیاب و مدل زبانی در یک زنجیرهٔ
          واحد \cite{lewis2020rag,gao2024ragsurvey}.
    \item[عاملیت و ابزارمندی] مدلی که خودش تصمیم می‌گیرد کِی جست‌وجو کند
          \cite{yao2023react,nakano2021webgpt}.
  \end{description}

  \begin{PTnote}
    این پژوهش در تلاقی دو خط دوم و سوم می‌نشیند: زنجیرهٔ روشن و سنجش‌پذیر، اما با
    اختیار \hl{جست‌وجوی زندهٔ وب} در زمان پاسخ‌دهی.
  \end{PTnote}
\end{frame}

%%%%%%%%%%%%%%%%%%%%%%%%%%%%%%%%%%%%%%%%%%%%%%%%%%%%%%%%%%%%%%%%%%%%%%%%%%%%%%%%
\begin{frame}{مقایسهٔ سامانه‌های موجود}
  \begin{center}
    \small
    \begin{tabular}{@{}lccc@{}}
      \toprule
      سامانه & بازیابی زندهٔ وب & ارجاع به منبع & تنظیم‌پذیری \\
      \midrule
      موتور جست‌وجوی کلیدواژه‌ای & دارد   & دارد   & ندارد \\
      دستیار گفت‌وگویی بدون ابزار & ندارد  & ندارد  & ندارد \\
      \en{Perplexity} و مانند آن  & دارد   & دارد   & جعبهٔ سیاه \\
      \midrule
      سامانهٔ پیشنهادی            & دارد   & دارد   & کامل \\
      \bottomrule
    \end{tabular}
  \end{center}

  \begin{itemize}
    \item ستون آخر تفاوت اصلی است: هر جزء زنجیره در سامانهٔ پیشنهادی
          \hl{جایگزین‌پذیر} و جداگانه \hl{سنجش‌پذیر} است.
  \end{itemize}
\end{frame}

%%%%%%%%%%%%%%%%%%%%%%%%%%%%%%%%%%%%%%%%%%%%%%%%%%%%%%%%%%%%%%%%%%%%%%%%%%%%%%%%
\begin{frame}{شکاف پژوهشی}
  \begin{PTcons}
    \begin{itemize}
      \item سامانه‌های تجاری نه ساختار درونی‌شان را آشکار می‌کنند و نه امکان
            سنجش جزء‌به‌جزء می‌دهند.
      \item پیاده‌سازی‌های پژوهشی معمولاً بر پیکرهٔ \hl{بسته و ایستا} کار می‌کنند،
            نه وب زنده.
    \end{itemize}
  \end{PTcons}

  \begin{PTresult}
    نیاز به سامانه‌ای \hl{سرتاسری، شفاف و سنجش‌پذیر} که روی وب زنده کار کند و
    هزینهٔ هر پرسش را نیز کنترل نماید.
  \end{PTresult}
\end{frame}

%%%%%%%%%%%%%%%%%%%%%%%%%%%%%%%%%%%%%%%%%%%%%%%%%%%%%%%%%%%%%%%%%%%%%%%%%%%%%%%%
%%  content/04-architecture.tex — فصل ۴: معماری پیشنهادی
%%
%%  نکته: متن فارسی داخل گره‌های TikZ همیشه در \PTrl{...} پیچیده می‌شود؛ توضیحش
%%  در بخش ۵ فایل theme/persian-thesis-boxes.sty آمده است.
%%%%%%%%%%%%%%%%%%%%%%%%%%%%%%%%%%%%%%%%%%%%%%%%%%%%%%%%%%%%%%%%%%%%%%%%%%%%%%%%
\section{معماری پیشنهادی}

%%%%%%%%%%%%%%%%%%%%%%%%%%%%%%%%%%%%%%%%%%%%%%%%%%%%%%%%%%%%%%%%%%%%%%%%%%%%%%%%
\begin{frame}{معماری کلی سامانه}{جریان پردازش از راست به چپ}
  \begin{center}
    \begin{tikzpicture}[x=1cm,y=1cm]
      % ---- پنج مرحلهٔ زنجیره -------------------------------------------------
      \node[PTnodemain]   (s1) at (11.40,0) {\PTrl{پرسش کاربر}};
      \node[PTnode]       (s2) at ( 8.55,0) {\PTrl{بازیابی زندهٔ وب}};
      \node[PTnodeaccent] (s3) at ( 5.70,0) {\PTrl{بازرتبه‌بندی دومرحله‌ای}};
      \node[PTnode]       (s4) at ( 2.85,0) {\PTrl{تولید پاسخ}};
      \node[PTnodemain]   (s5) at ( 0.00,0) {\PTrl{پاسخ مستند}};

      % ---- پیکان‌ها ----------------------------------------------------------
      \draw[PTflow] (s1.west) -- (s2.east);
      \draw[PTflow] (s2.west) -- (s3.east);
      \draw[PTflow] (s3.west) -- (s4.east);
      \draw[PTflow] (s4.west) -- (s5.east);

      % ---- فناوری هر مرحله --------------------------------------------------
      \node[PTcapt,anchor=north] at ( 8.55,-0.74) {\en{Tavily Search API}};
      % «\\» بیرون از \en می‌ماند: شکستن خط کار خودِ گره است، نه کار یک بازهٔ چپ‌به‌راست.
      \node[PTcapt,anchor=north] at ( 5.70,-0.74) {\en{Cohere Rerank}\\\en{+ Jina Reranker}};
      \node[PTcapt,anchor=north] at ( 2.85,-0.74) {\en{Gemini 2.5 Flash}};

      % ---- برجسته‌سازی مشارکت اصلی ------------------------------------------
      \draw[PTaccent,line width=.7pt]
        (4.45,0.72) -- (4.45,0.96) -- (6.95,0.96) -- (6.95,0.72);
      \node[text=PTaccentdark,font=\tiny,anchor=south] at (5.70,0.99)
        {\PTrl{مشارکت اصلی این پژوهش}};
    \end{tikzpicture}
  \end{center}

  \begin{itemize}
    \item هر مرحله یک \term{واسط روشن} دارد، پس می‌توان جایگزینش کرد و اثرش را
          جداگانه سنجید.
    \item بازرتبه‌بندی، بافتار ورودی مدل را از ده‌ها سند به \hl{سه سند} کوتاه
          می‌کند.
  \end{itemize}
\end{frame}

%%%%%%%%%%%%%%%%%%%%%%%%%%%%%%%%%%%%%%%%%%%%%%%%%%%%%%%%%%%%%%%%%%%%%%%%%%%%%%%%
\begin{frame}{شرح مرحله‌به‌مرحله}
  \begin{description}
    \item[بازنویسی پرسش] پرسش کاربر به چند پرس‌وجوی جست‌وجو گسترده می‌شود تا
          یادآوری بالا برود.
    \item[بازیابی] جست‌وجوی زندهٔ وب و برداشت متن صفحات، همراه با نگه‌داشتن
          \term{نشانی منبع} برای هر قطعه.
    \item[بازرتبه‌بندی ۱] مدل \en{Cohere Rerank} فهرست خام را به ده سند برتر
          کوچک می‌کند.
    \item[بازرتبه‌بندی ۲] مدل \en{Jina Reranker} همان ده سند را به \hl{سه سند}
          نهایی می‌رساند.
    \item[تولید] مدل زبانی تنها با این سه سند پاسخ می‌سازد و در متن به آن‌ها
          ارجاع می‌دهد.
  \end{description}

  \begin{PTnote}
    دو مرحله‌ای کردن بازرتبه‌بندی، به‌جای یک مرحلهٔ بزرگ، دقت را حفظ می‌کند و
    هزینهٔ فراخوانی را پایین می‌آورد.
  \end{PTnote}
\end{frame}

%%%%%%%%%%%%%%%%%%%%%%%%%%%%%%%%%%%%%%%%%%%%%%%%%%%%%%%%%%%%%%%%%%%%%%%%%%%%%%%%
\begin{frame}{پشتهٔ پیاده‌سازی}
  \begin{center}
    \small
    \begin{tabular}{@{}ll@{}}
      \toprule
      لایه & فناوری \\
      \midrule
      کارساز و واسط برنامه‌نویسی & \en{Python 3.12} + \en{FastAPI} \\
      پایگاه داده                & \en{SQLite} + \en{SQLAlchemy} \\
      واسط کاربری                & \en{HTML5} + \en{Tailwind CSS v4} \\
      مدل زبانی                  & \en{Gemini 2.5 Flash} \\
      جست‌وجوی وب                & \en{Tavily Search API} \\
      بازرتبه‌بندی               & \en{Cohere Rerank}، \en{Jina Reranker} \\
      \bottomrule
    \end{tabular}
  \end{center}

  \begin{itemize}
    \item همهٔ سرویس‌های بیرونی پشت یک \term{لایهٔ واسط} قرار دارند؛ تعویض هر یک
          فقط یک کلاس را تغییر می‌دهد.
  \end{itemize}
\end{frame}

%%%%%%%%%%%%%%%%%%%%%%%%%%%%%%%%%%%%%%%%%%%%%%%%%%%%%%%%%%%%%%%%%%%%%%%%%%%%%%%%
%%  اسلاید کد
%%  ۱) هر اسلایدی که PTcode دارد باید [fragile] بگیرد.
%%  ۲) داخل PTcode فقط نویسهٔ اَسکی بنویسید. قلم تک‌عرض پیش‌فرض نشانهٔ فارسی ندارد
%%     و listings هم با نویسه‌های چندبایتی خوب کنار نمی‌آید؛ پس توضیحات کد را
%%     انگلیسی بنویسید و شرح فارسی را در فهرست زیر اسلاید بیاورید.
%%%%%%%%%%%%%%%%%%%%%%%%%%%%%%%%%%%%%%%%%%%%%%%%%%%%%%%%%%%%%%%%%%%%%%%%%%%%%%%%
\begin{frame}[fragile]{هستهٔ زنجیره}{ساده‌سازی‌شده برای ارائه}
\begin{PTcode}[language=Python]
async def answer(question: str) -> Answer:
    queries   = rewrite(question)              # query expansion
    documents = await search_web(queries)      # Tavily
    top10     = cohere_rerank(question, documents, k=10)
    top3      = jina_rerank(question, top10,   k=3)
    return generate(question, context=top3)    # Gemini 2.5 Flash
\end{PTcode}

  \begin{itemize}
    \item کل زنجیره در یک تابع ناهم‌زمان جمع می‌شود؛ هر گام یک وابستگی
          جایگزین‌پذیر است.
    \item دو فراخوانی \en{rerank} پشت سر هم می‌آیند: نخست ده سند، سپس سه سند.
  \end{itemize}
\end{frame}

%%%%%%%%%%%%%%%%%%%%%%%%%%%%%%%%%%%%%%%%%%%%%%%%%%%%%%%%%%%%%%%%%%%%%%%%%%%%%%%%
%%  content/05-evaluation.tex — فصل ۵: ارزیابی و نتایج
%%
%%  نمودار ستونی این فصل با TikZ خالص کشیده شده است تا به pgfplots نیازی نباشد.
%%  ترفند مقیاس: با y=0.032cm هر واحد محور عمودی برابر یک درصد است، پس مقدارها را
%%  می‌توان مستقیم به‌صورت درصد نوشت و ۱۰۰ درصد ۳٫۲ سانتی‌متر بلند می‌شود.
%%%%%%%%%%%%%%%%%%%%%%%%%%%%%%%%%%%%%%%%%%%%%%%%%%%%%%%%%%%%%%%%%%%%%%%%%%%%%%%%
\section{ارزیابی و نتایج}

%%%%%%%%%%%%%%%%%%%%%%%%%%%%%%%%%%%%%%%%%%%%%%%%%%%%%%%%%%%%%%%%%%%%%%%%%%%%%%%%
\begin{frame}{روش ارزیابی}
  \begin{description}
    \item[مجموعهٔ آزمون] \hl{۴۰ پرسش} در چهار دستهٔ رخداد تازه، پرسش چندبخشی،
          پرسش عددی و پرسش تعریفی.
    \item[سنجه‌ها] درستی پاسخ (داوری انسانی)، \en{Precision@3} و \en{NDCG@3}
          برای مرحلهٔ بازیابی.
    \item[مبنای مقایسه] یک سامانهٔ تجاری پرسش و پاسخ مبتنی بر وب، با همان
          پرسش‌ها و در همان بازهٔ زمانی.
    \item[مطالعهٔ کاهشی] اجرای زنجیره یک‌بار \term{بدون} بازرتبه‌بندی و یک‌بار
          \term{با} آن، برای جدا کردن سهم این مرحله.
  \end{description}

  \begin{PTnote}
    داوری روی درستی پاسخ به‌صورت دوسویه و بدون آگاهی از نام سامانه انجام شد تا
    سوگیری داور کنترل شود.
  \end{PTnote}
\end{frame}

%%%%%%%%%%%%%%%%%%%%%%%%%%%%%%%%%%%%%%%%%%%%%%%%%%%%%%%%%%%%%%%%%%%%%%%%%%%%%%%%
\begin{frame}{مقایسه با سامانهٔ مرجع}{۴۰ پرسش، داوری انسانی}
  \begin{center}
    \small
    \begin{tabular}{@{}lcc@{}}
      \toprule
      سنجه & سامانهٔ پیشنهادی & سامانهٔ مرجع \\
      \midrule
      درستی پاسخ            & \hl{۹۵٪}   & ۷۷٫۵٪ \\
      پاسخ همراه ارجاع      & ۱۰۰٪       & ۴۵٪   \\
      میانگین اسناد بافتار  & \hl{۳}     & ۱۰ تا ۲۰ \\
      \bottomrule
    \end{tabular}
  \end{center}

  \begin{PTresult}
    سامانهٔ پیشنهادی با بافتاری \hl{کوتاه‌تر}، درستی بالاتری به دست می‌آورد؛ پس
    بهبود از \term{کیفیت بازیابی} می‌آید، نه از حجم بافتار.
  \end{PTresult}
\end{frame}

%%%%%%%%%%%%%%%%%%%%%%%%%%%%%%%%%%%%%%%%%%%%%%%%%%%%%%%%%%%%%%%%%%%%%%%%%%%%%%%%
\begin{frame}{اثر بازرتبه‌بندی بر کیفیت بازیابی}
  \begin{center}
    \begin{tikzpicture}[
        x = 1cm, y = 0.032cm,
        PTbarA/.style = {fill=PTaccentdark, draw=PTaccentdark, line width=.4pt},
        PTbarB/.style = {fill=PTline,       draw=PTmuted,      line width=.4pt},
        PTgrid/.style = {draw=PTline, line width=.4pt,
                         dash pattern=on 1pt off 1.4pt},
        PTval/.style  = {font=\tiny, text=PTink,  anchor=south, inner sep=1pt},
        PTswatch/.style = {inner sep=0pt, minimum width=0.34cm,
                           minimum height=0.20cm},
      ]
      % ---- شبکه و محورها ------------------------------------------------------
      \draw[PTgrid] (1.6, 25) -- (11.4, 25);
      \draw[PTgrid] (1.6, 50) -- (11.4, 50);
      \draw[PTgrid] (1.6, 75) -- (11.4, 75);
      \draw[PTgrid] (1.6,100) -- (11.4,100);
      \draw[draw=PTmuted,line width=.6pt] (11.4,0) -- (11.4,106);
      \draw[draw=PTmuted,line width=.6pt] ( 1.6,0) -- (11.4,  0);

      % ---- برچسب محور عمودی: سمت راست، چون خواندن از راست آغاز می‌شود ---------
      \node[PTcapt,anchor=west] at (11.55,  0) {\PTrl{۰}};
      \node[PTcapt,anchor=west] at (11.55, 25) {\PTrl{۲۵}};
      \node[PTcapt,anchor=west] at (11.55, 50) {\PTrl{۵۰}};
      \node[PTcapt,anchor=west] at (11.55, 75) {\PTrl{۷۵}};
      \node[PTcapt,anchor=west] at (11.55,100) {\PTrl{۱۰۰}};

      % ---- گروه ۱ (راست): Precision@3 ---------------------------------------
      \fill[PTbarB] (7.48,0) rectangle (8.53,67.7);
      \fill[PTbarA] (8.67,0) rectangle (9.72,93.3);
      \node[PTval] at (8.005,67.7) {\PTrl{۶۷٫۷٪}};
      \node[PTval] at (9.195,93.3) {\PTrl{۹۳٫۳٪}};
      \node[PTcapt,anchor=north] at (8.6,-4) {\en{Precision@3}};

      % ---- گروه ۲ (چپ): NDCG@3 ----------------------------------------------
      \fill[PTbarB] (3.28,0) rectangle (4.33,82.4);
      \fill[PTbarA] (4.47,0) rectangle (5.52,94.7);
      \node[PTval] at (3.805,82.4) {\PTrl{۸۲٫۴٪}};
      \node[PTval] at (4.995,94.7) {\PTrl{۹۴٫۷٪}};
      \node[PTcapt,anchor=north] at (4.4,-4) {\en{NDCG@3}};

      % ---- راهنما (از راست به چپ) --------------------------------------------
      \node[PTswatch,fill=PTaccentdark] at (9.30,116) {};
      \node[PTcapt,anchor=east]         at (9.08,116) {\PTrl{با بازرتبه‌بندی}};
      \node[PTswatch,fill=PTline,draw=PTmuted,line width=.4pt]
                                        at (6.10,116) {};
      \node[PTcapt,anchor=east]         at (5.88,116) {\PTrl{بدون بازرتبه‌بندی}};
    \end{tikzpicture}
  \end{center}

  \begin{itemize}
    \item \en{Precision@3} حدود \hl{۲۵٫۶ واحد درصد} بالا می‌رود؛ یعنی سه سند
          نهایی تقریباً همیشه مرتبط‌اند.
    \item رشد کمترِ \en{NDCG@3} نشان می‌دهد بازیابی پایه اسناد درست را
          \term{پیدا} می‌کرده، اما \term{مرتب} نمی‌کرده است.
  \end{itemize}
\end{frame}

%%%%%%%%%%%%%%%%%%%%%%%%%%%%%%%%%%%%%%%%%%%%%%%%%%%%%%%%%%%%%%%%%%%%%%%%%%%%%%%%
\begin{frame}{تحلیل خطا}
  \begin{PTcons}
    \begin{itemize}
      \item دو پرسش نادرست، هر دو از دستهٔ \term{عددی}: منبع درست بازیابی شد اما
            عدد در جدول صفحه بود و از متن برداشت نشد.
      \item پرسش‌هایی که پاسخشان در وب پراکنده است، به بیش از سه سند نیاز دارند.
    \end{itemize}
  \end{PTcons}

  \begin{itemize}
    \item تأخیر میانگین هر پرسش \hl{۴٫۲ ثانیه} است که بیشترش صرف بازیابی زندهٔ
          وب می‌شود، نه بازرتبه‌بندی.
    \item هزینهٔ هر پرسش نسبت به سپردن همهٔ اسناد به مدل، حدود \hl{یک‌سوم} است.
  \end{itemize}
\end{frame}

%%%%%%%%%%%%%%%%%%%%%%%%%%%%%%%%%%%%%%%%%%%%%%%%%%%%%%%%%%%%%%%%%%%%%%%%%%%%%%%%
%%  content/06-conclusion.tex — فصل ۶: نتیجه‌گیری و کارهای آینده
%%%%%%%%%%%%%%%%%%%%%%%%%%%%%%%%%%%%%%%%%%%%%%%%%%%%%%%%%%%%%%%%%%%%%%%%%%%%%%%%
\section{نتیجه‌گیری و کارهای آینده}

%%%%%%%%%%%%%%%%%%%%%%%%%%%%%%%%%%%%%%%%%%%%%%%%%%%%%%%%%%%%%%%%%%%%%%%%%%%%%%%%
\begin{frame}{جمع‌بندی}
  \begin{itemize}
    \item سامانه‌ای ساخته شد که پرسش زبان طبیعی را می‌گیرد و پاسخی
          \hl{مستند} با ارجاع به منبع بازمی‌گرداند.
    \item زنجیرهٔ \term{بازیابی، بازرتبه‌بندی و تولید} به اجزای جایگزین‌پذیر شکسته
          شد، پس هر مرحله جداگانه سنجیدنی است.
    \item \term{بازرتبه‌بندی دومرحله‌ای} بافتار ورودی مدل را به سه سند کوتاه کرد و
          در همان حال دقت بازیابی را بالا برد.
  \end{itemize}

  \begin{PTresult}
    درستی پاسخ \hl{۹۵٪} در برابر \hl{۷۷٫۵٪} سامانهٔ مرجع، با بافتاری کوتاه‌تر و
    هزینهٔ کمتر برای هر پرسش.
  \end{PTresult}
\end{frame}

%%%%%%%%%%%%%%%%%%%%%%%%%%%%%%%%%%%%%%%%%%%%%%%%%%%%%%%%%%%%%%%%%%%%%%%%%%%%%%%%
\begin{frame}{محدودیت‌ها}
  \begin{PTcons}
    \begin{itemize}
      \item وابستگی به سرویس‌های بیرونی: هزینه، سهمیهٔ فراخوانی و دسترس‌پذیری
            خارج از کنترل سامانه است.
      \item سه سند نهایی برای پرسش‌هایی که پاسخشان در وب پراکنده است کم است.
      \item مجموعهٔ آزمون \hl{۴۰ پرسشی} کوچک است و نتیجه را باید مقدماتی دانست.
    \end{itemize}
  \end{PTcons}

  \begin{itemize}
    \item هیچ‌یک از این محدودیت‌ها به \term{معماری} گره نخورده است؛ هر سه با
          تعویض یا تنظیم یک جزء برطرف می‌شوند.
  \end{itemize}
\end{frame}

%%%%%%%%%%%%%%%%%%%%%%%%%%%%%%%%%%%%%%%%%%%%%%%%%%%%%%%%%%%%%%%%%%%%%%%%%%%%%%%%
\begin{frame}{کارهای آینده}
  \begin{enumerate}
    \item \term{تعداد سند پویا}: انتخاب \(k\) بر پایهٔ امتیاز بازرتبه‌بند، به‌جای
          عدد ثابت سه.
    \item \term{بازرتبه‌بند بومی}: ریزتنظیم یک مدل رمزگذار متقابل روی دادهٔ فارسی
          برای حذف وابستگی به سرویس بیرونی.
    \item \term{حافظهٔ گفت‌وگو}: پشتیبانی از پرسش‌های پیاپی و ارجاع‌های درون‌متنی.
    \item \term{ارزیابی گسترده‌تر}: مجموعهٔ چند‌صد پرسشی با داوری چنددادوری و
          گزارش بازهٔ اطمینان.
  \end{enumerate}

  \begin{PTnote}
    ترتیب بالا بر پایهٔ نسبت \hl{اثر به هزینه} چیده شده است؛ گام نخست کمترین
    تغییر در کد و بیشترین اثر را دارد.
  \end{PTnote}
\end{frame}


% ---- مراجع ----------------------------------------------------------------
%%%%%%%%%%%%%%%%%%%%%%%%%%%%%%%%%%%%%%%%%%%%%%%%%%%%%%%%%%%%%%%%%%%%%%%%%%%%%%%%
%%  content/07-references.tex — مراجع
%%
%%  چرا thebibliography دستی و نه biblatex؟
%%  برای اینکه قالب با یک بار اجرای xelatex هم کامل در بیاید و به bibtex/biber
%%  نیازی نداشته باشد.  اگر biblatex می‌خواهید، همین قاب را با \printbibliography
%%  جایگزین کنید و در main.tex این دو خط را پیش از config/fonts.tex بیفزایید:
%%
%%      \usepackage[backend=biber,style=numeric-comp,sorting=none]{biblatex}
%%      \addbibresource{refs/references.bib}
%%
%%  فهرست انگلیسی در محیط PTltr پیچیده شده است: یک بلوکِ چپ‌به‌راستِ بندیْ که
%%  خطهایش می‌شکنند و می‌تواند از یک اسلاید به اسلاید بعد سرریز کند.
%%%%%%%%%%%%%%%%%%%%%%%%%%%%%%%%%%%%%%%%%%%%%%%%%%%%%%%%%%%%%%%%%%%%%%%%%%%%%%%%
\begin{frame}[allowframebreaks]{\PTlabel{references}}
  \begin{PTltr}
    \scriptsize
    \begin{thebibliography}{99}

      \bibitem{lewis2020rag}
      P.~Lewis \textit{et al.}, ``Retrieval-Augmented Generation for
      Knowledge-Intensive NLP Tasks,'' in \textit{Advances in Neural Information
      Processing Systems (NeurIPS)}, 2020.

      \bibitem{karpukhin2020dpr}
      V.~Karpukhin \textit{et al.}, ``Dense Passage Retrieval for Open-Domain
      Question Answering,'' in \textit{Proc. EMNLP}, 2020, pp.~6769--6781.

      \bibitem{gao2024ragsurvey}
      Y.~Gao \textit{et al.}, ``Retrieval-Augmented Generation for Large Language
      Models: A Survey,'' \textit{arXiv:2312.10997}, 2024.

      \bibitem{vaswani2017}
      A.~Vaswani \textit{et al.}, ``Attention Is All You Need,'' in
      \textit{Advances in Neural Information Processing Systems (NeurIPS)}, 2017.

      \bibitem{devlin2019bert}
      J.~Devlin, M.-W.~Chang, K.~Lee, and K.~Toutanova, ``BERT: Pre-training of
      Deep Bidirectional Transformers for Language Understanding,'' in
      \textit{Proc. NAACL-HLT}, 2019, pp.~4171--4186.

      \bibitem{reimers2019sbert}
      N.~Reimers and I.~Gurevych, ``Sentence-BERT: Sentence Embeddings using
      Siamese BERT-Networks,'' in \textit{Proc. EMNLP-IJCNLP}, 2019.

      \bibitem{robertson2009bm25}
      S.~Robertson and H.~Zaragoza, ``The Probabilistic Relevance Framework:
      BM25 and Beyond,'' \textit{Foundations and Trends in Information
      Retrieval}, vol.~3, no.~4, pp.~333--389, 2009.

      \bibitem{nogueira2019passage}
      R.~Nogueira and K.~Cho, ``Passage Re-ranking with BERT,''
      \textit{arXiv:1901.04085}, 2019.

      \bibitem{khattab2020colbert}
      O.~Khattab and M.~Zaharia, ``ColBERT: Efficient and Effective Passage
      Search via Contextualized Late Interaction over BERT,'' in
      \textit{Proc. SIGIR}, 2020, pp.~39--48.

      \framebreak

      \bibitem{izacard2021fid}
      G.~Izacard and E.~Grave, ``Leveraging Passage Retrieval with Generative
      Models for Open Domain Question Answering,'' in \textit{Proc. EACL}, 2021,
      pp.~874--880.

      \bibitem{yao2023react}
      S.~Yao \textit{et al.}, ``ReAct: Synergizing Reasoning and Acting in
      Language Models,'' in \textit{Proc. ICLR}, 2023.

      \bibitem{nakano2021webgpt}
      R.~Nakano \textit{et al.}, ``WebGPT: Browser-assisted Question-answering
      with Human Feedback,'' \textit{arXiv:2112.09332}, 2021.

      \bibitem{asai2024selfrag}
      A.~Asai, Z.~Wu, Y.~Wang, A.~Sil, and H.~Hajishirzi, ``Self-RAG: Learning
      to Retrieve, Generate, and Critique through Self-Reflection,'' in
      \textit{Proc. ICLR}, 2024.

      \bibitem{es2024ragas}
      S.~Es, J.~James, L.~Espinosa-Anke, and S.~Schockaert, ``RAGAS: Automated
      Evaluation of Retrieval Augmented Generation,'' in \textit{Proc. EACL:
      System Demonstrations}, 2024, pp.~150--158.

      \bibitem{jarvelin2002ndcg}
      K.~J\"arvelin and J.~Kek\"al\"ainen, ``Cumulated Gain-based Evaluation of
      IR Techniques,'' \textit{ACM Transactions on Information Systems},
      vol.~20, no.~4, pp.~422--446, 2002.

    \end{thebibliography}
  \end{PTltr}
\end{frame}


% ---- اسلاید پایانی --------------------------------------------------------
\PTclosingframe

% ---- اسلایدهای پشتیبان (پس از اسلاید تشکر، برای پاسخ به پرسش‌ها) ----------
%%%%%%%%%%%%%%%%%%%%%%%%%%%%%%%%%%%%%%%%%%%%%%%%%%%%%%%%%%%%%%%%%%%%%%%%%%%%%%%%
%%  content/08-appendix.tex — اسلایدهای پشتیبان
%%
%%  این فایل *پس از* \PTclosingframe در main.tex می‌آید: اسلایدهایی که در ارائهٔ
%%  اصلی نشان داده نمی‌شوند و فقط برای پاسخ به پرسش داوران آماده‌اند.
%%
%%  همهٔ قاب‌ها [noframenumbering] دارند، پس شمارهٔ اسلاید و نوار پیشرفت دست
%%  نمی‌خورند و «۱۸ از ۱۸» در پایان ارائه همان ۱۸ می‌ماند.
%%%%%%%%%%%%%%%%%%%%%%%%%%%%%%%%%%%%%%%%%%%%%%%%%%%%%%%%%%%%%%%%%%%%%%%%%%%%%%%%
\appendix

\section{پیوست: اسلایدهای پشتیبان}

%%%%%%%%%%%%%%%%%%%%%%%%%%%%%%%%%%%%%%%%%%%%%%%%%%%%%%%%%%%%%%%%%%%%%%%%%%%%%%%%
\begin{frame}[noframenumbering]{چرا دو بازرتبه‌بند و نه یکی؟}
  \begin{center}
    \small
    \begin{tabular}{@{}lccc@{}}
      \toprule
      پیکربندی & \en{Precision@3} & هزینهٔ نسبی & تأخیر \\
      \midrule
      بدون بازرتبه‌بندی        & ۶۷٫۷٪    & ۱٫۰ & ۳٫۱ ثانیه \\
      فقط \en{Cohere}          & ۸۸٫۱٪    & ۱٫۲ & ۳٫۶ ثانیه \\
      فقط \en{Jina}            & ۸۶٫۴٪    & ۱٫۵ & ۳٫۹ ثانیه \\
      \en{Cohere} + \en{Jina}  & \hl{۹۳٫۳٪} & ۱٫۳ & ۴٫۲ ثانیه \\
      \bottomrule
    \end{tabular}
  \end{center}

  \begin{itemize}
    \item مرحلهٔ نخست \term{ارزان و گسترده} است و فهرست خام را کوچک می‌کند؛ مرحلهٔ
          دوم \term{گران و باریک} است و فقط روی ده سند اجرا می‌شود.
    \item همین ترتیب است که هزینهٔ کل را از اجرای یک‌مرحله‌ای گران‌تر پایین‌تر
          نگه می‌دارد.
  \end{itemize}
\end{frame}

%%%%%%%%%%%%%%%%%%%%%%%%%%%%%%%%%%%%%%%%%%%%%%%%%%%%%%%%%%%%%%%%%%%%%%%%%%%%%%%%
\begin{frame}[noframenumbering]{فراپارامترها و تنظیمات}
  \begin{center}
    \small
    \begin{tabular}{@{}ll@{}}
      \toprule
      پارامتر & مقدار \\
      \midrule
      تعداد پرس‌وجوی بازنویسی‌شده          & ۳ \\
      تعداد نتیجهٔ هر پرس‌وجو              & ۱۰ \\
      خروجی بازرتبه‌بند نخست               & ۱۰ سند \\
      خروجی بازرتبه‌بند دوم                & ۳ سند \\
      بیشینهٔ طول هر قطعهٔ متن             & ۱۰۰۰ نویسه \\
      دمای مدل زبانی                       & ۰٫۲ \\
      \bottomrule
    \end{tabular}
  \end{center}

  \begin{PTnote}
    دمای پایین انتخاب شده است تا مدل از سه سند داده‌شده فراتر نرود؛ همین تنظیم
    نرخ توهم‌زایی را در آزمون‌ها به صفر رساند.
  \end{PTnote}
\end{frame}

%%%%%%%%%%%%%%%%%%%%%%%%%%%%%%%%%%%%%%%%%%%%%%%%%%%%%%%%%%%%%%%%%%%%%%%%%%%%%%%%
\begin{frame}[noframenumbering]{نمونهٔ پاسخ سامانه}
  \begin{PTexample}
    \textbf{پرسش:} آخرین نسخهٔ پایدار پایتون چیست و چه زمانی منتشر شد؟

    \medskip
    \textbf{پاسخ:} نسخهٔ پایدار کنونی \en{Python 3.12} است \soft{[۱]} که در
    مهر ۱۴۰۲ منتشر شد \soft{[۲]}. مهم‌ترین تغییرش بهبود پیام‌های خطا و
    کارایی مفسر است \soft{[۱]}.
  \end{PTexample}

  \begin{itemize}
    \item هر جملهٔ پاسخ به شمارهٔ سندی گره خورده است که از آن برداشت شده؛ داور
          می‌تواند مستقیم به منبع برود.
    \item شماره‌ها به نشانی کامل صفحه در پانویس پاسخ نگاشته می‌شوند.
  \end{itemize}
\end{frame}

%%%%%%%%%%%%%%%%%%%%%%%%%%%%%%%%%%%%%%%%%%%%%%%%%%%%%%%%%%%%%%%%%%%%%%%%%%%%%%%%
\begin{frame}[noframenumbering]{ساختار دستهٔ پرسش‌های آزمون}
  \begin{center}
    \small
    \begin{tabular}{@{}lcc@{}}
      \toprule
      دستهٔ پرسش & تعداد & درستی \\
      \midrule
      رخداد تازه     & ۱۲ & ۱۰۰٪ \\
      چندبخشی        & ۱۰ & ۹۰٪  \\
      عددی           & ۸  & ۷۵٪  \\
      تعریفی         & ۱۰ & ۱۰۰٪ \\
      \midrule
      مجموع          & ۴۰ & \hl{۹۵٪} \\
      \bottomrule
    \end{tabular}
  \end{center}

  \begin{itemize}
    \item ضعیف‌ترین دسته، پرسش‌های \term{عددی} است؛ ریشهٔ خطا برداشت‌نشدن عدد از
          جدول‌های صفحهٔ وب است، نه بازیابی نادرست.
  \end{itemize}
\end{frame}


\end{document}
